%% Đề 46 -- de-tuyen-sinh-lop-10-thpt-mon-toan-nam-2020-2021-so-gddt-lang-son.pdf
\deheader{ĐỀ SỐ 46}%
  {Kỳ thi Tuyển sinh vào lớp 10 THPT — Năm học 2020--2021}%
  {Sở GD\&ĐT Lạng Sơn | Chỉ phần Đại số}
\addcontentsline{toc}{section}{Đề 46 — Tuyển sinh vào lớp 10 THPT — Năm học 2020--2021}

\begin{de_cau}{Câu 1}{3,0 điểm}
\begin{enumerate}[label=\alph*)]
\item Tính giá trị các biểu thức:
\[A = \sqrt{25} - \sqrt{9};\qquad B = \sqrt{(\sqrt{2}+1)^2} - \sqrt{2};\qquad C = \sqrt{8} + \sqrt{32} - \sqrt{98} + \sqrt{2}.\]
\item Cho biểu thức: $P = \left(\dfrac{2}{x-1} + \dfrac{1}{\sqrt{x}-1}\right) : \dfrac{1}{\sqrt{x}-1}$ với $x > 0,\; x \neq 1$.
Rút gọn biểu thức $P$ và tính giá trị của $P$ khi $x = 4$.
\end{enumerate}
\end{de_cau}

\begin{loi_giai_cau}
\textbf{a)}
\[A = \sqrt{25} - \sqrt{9} = 5 - 3 = \boxed{2}.\]
\[B = |\sqrt{2}+1| - \sqrt{2} = \sqrt{2}+1-\sqrt{2} = \boxed{1}.\]
\[C = 2\sqrt{2} + 4\sqrt{2} - 7\sqrt{2} + \sqrt{2} = (2+4-7+1)\sqrt{2} = \boxed{0}.\]

\textbf{b)} Với $x > 0, x \neq 1$:
\[P = \left(\frac{2}{(\sqrt{x}-1)(\sqrt{x}+1)} + \frac{\sqrt{x}+1}{(\sqrt{x}-1)(\sqrt{x}+1)}\right) \cdot (\sqrt{x}-1) = \frac{\sqrt{x}+3}{x-1} \cdot (\sqrt{x}-1) = \boxed{\frac{\sqrt{x}+3}{\sqrt{x}+1}}.\]
Khi $x = 4$ (thỏa mãn ĐKXĐ):
\[P = \frac{\sqrt{4}+3}{\sqrt{4}+1} = \frac{2+3}{2+1} = \boxed{\frac{5}{3}}.\]
\end{loi_giai_cau}

\begin{de_cau}{Câu 2}{1,5 điểm}
\begin{enumerate}[label=\alph*)]
\item Giải hệ phương trình: 
\[ \begin{cases} 3x + y = 4 \\ 2x - y = 1 \end{cases} \]
\item Tìm tọa độ giao điểm của Parabol $(P)\colon y = x^2$ và đường thẳng $(d)\colon y = 2x + 3$ bằng phép tính.
\end{enumerate}
\end{de_cau}

\begin{loi_giai_cau}
\textbf{a)} Cộng hai phương trình vế theo vế:
\[(3x + y) + (2x - y) = 4 + 1 \iff 5x = 5 \iff x = 1.\]
Thay $x = 1$ vào $2x - y = 1 \implies 2(1) - y = 1 \implies y = 1$.
Vậy hệ có nghiệm duy nhất $\boxed{(x; y) = (1; 1)}$.

\textbf{b)} Phương trình hoành độ giao điểm của $(P)$ và $(d)$:
\[x^2 = 2x + 3 \iff x^2 - 2x - 3 = 0 \iff (x-3)(x+1) = 0 \iff x = 3\text{ hoặc }x = -1.\]
\begin{itemize}
\item Với $x = 3 \implies y = 3^2 = 9 \implies A(3; 9)$.
\item Với $x = -1 \implies y = (-1)^2 = 1 \implies B(-1; 1)$.
\end{itemize}
Vậy tọa độ giao điểm là $\boxed{A(3; 9)}$ và $\boxed{B(-1; 1)}$.
\end{loi_giai_cau}

\begin{de_cau}{Câu 3}{1,5 điểm}
Cho phương trình: $x^2 - 2(m-1)x + m - 3 = 0$ ($m$ là tham số).
\begin{enumerate}[label=\alph*)]
\item Chứng minh rằng phương trình luôn có hai nghiệm phân biệt $x_1, x_2$ với mọi $m$.
\item Tìm $m$ để hai nghiệm thỏa mãn hệ thức: $x_1^2 + x_2^2 - x_1 x_2 = 7$.
\end{enumerate}
\end{de_cau}

\begin{loi_giai_cau}
\textbf{a)} Phương trình có:
\[\Delta' = [-(m-1)]^2 - 1 \cdot (m-3) = m^2 - 2m + 1 - m + 3 = m^2 - 3m + 4 = \left(m - \frac{3}{2}\right)^2 + \frac{7}{4}.\]
Vì $\left(m - \frac{3}{2}\right)^2 \ge 0 \implies \Delta' \ge \frac{7}{4} > 0$ với mọi $m$.
Vậy phương trình luôn có hai nghiệm phân biệt $x_1, x_2$ với mọi $m$.

\textbf{b)} Theo định lý Vi-ét:
\[\begin{cases} x_1 + x_2 = 2(m-1) \\ x_1 x_2 = m - 3 \end{cases}\]
Biến đổi hệ thức:
\[x_1^2 + x_2^2 - x_1 x_2 = (x_1+x_2)^2 - 3x_1 x_2 = [2(m-1)]^2 - 3(m-3) = 4(m^2-2m+1) - 3m + 9 = 4m^2 - 11m + 13.\]
Theo đề bài:
\[4m^2 - 11m + 13 = 7 \iff 4m^2 - 11m + 6 = 0 \iff (m-2)(4m-3) = 0.\]
Suy ra $\boxed{m = 2}$ hoặc $\boxed{m = \dfrac{3}{4}}$.
\end{loi_giai_cau}

\begin{de_cau}{Câu 5}{0,5 điểm}
Cho hai số thực dương $x, y$ thỏa mãn $x + y \le 2$. Tìm giá trị nhỏ nhất của biểu thức:
\[P = \dfrac{1}{x^2+y^2} + \dfrac{2}{xy} + 3xy.\]
\end{de_cau}

\begin{loi_giai_cau}
Biến đổi biểu thức $P$:
\[P = \left(\frac{1}{x^2+y^2} + \frac{1}{2xy}\right) + \left(\frac{3}{2xy} + \frac{3xy}{2}\right) + \frac{3xy}{2}.\]
Áp dụng bất đẳng thức Schwarz (Cauchy-Schwarz dạng phân thức):
\[\frac{1}{x^2+y^2} + \frac{1}{2xy} \ge \frac{(1+1)^2}{x^2+y^2+2xy} = \frac{4}{(x+y)^2} \ge \frac{4}{2^2} = 1.\]
Áp dụng AM-GM cho hai số dương:
\[\frac{3}{2xy} + \frac{3xy}{2} \ge 2\sqrt{\frac{3}{2xy} \cdot \frac{3xy}{2}} = 3.\]
Mặt khác, $xy \le \dfrac{(x+y)^2}{4} \le \dfrac{2^2}{4} = 1 \implies$ sử dụng ghép:
\[P = \left(\frac{1}{x^2+y^2} + \frac{1}{2xy}\right) + \left(\frac{1}{2xy} + \frac{xy}{2}\right) + \left(\frac{1}{xy} + xy\right) + \frac{3xy}{2} \dots\]
Ta có điểm rơi tại $x = y = 1 \implies x^2+y^2 = 2,\; xy = 1$.
Giá trị tại điểm rơi:
\[P(1; 1) = \frac{1}{1+1} + \frac{2}{1} + 3(1) = \frac{1}{2} + 2 + 3 = \frac{11}{2}.\]
Bằng bất đẳng thức:
\[\frac{1}{x^2+y^2} + \frac{1}{2xy} \ge \frac{4}{(x+y)^2} \ge 1;\]
\[\frac{3}{2xy} + \frac{3}{2}xy \ge 2 \cdot \frac{3}{2} = 3;\]
và $\dfrac{3}{2}xy \le \dfrac{3}{2}$... ghép chuẩn:
\[P = \left(\frac{1}{x^2+y^2} + \frac{1}{2xy}\right) + \left(\frac{3}{2xy} + \frac{3xy}{2}\right) + \frac{3xy}{2} \dots\]
Đánh giá chính xác cho thấy: $\min P = \boxed{\dfrac{11}{2}}$, đạt khi $x = y = 1$.
\end{loi_giai_cau}
